% =============================================================================
%  Exosuit Meeting Notes — April 2026
%  Technical analysis of design decisions discussed with Chris Kalogroulis
% =============================================================================

\documentclass[11pt, a4paper]{article}

\usepackage{amsmath, amssymb, bm}
\usepackage{geometry}
\usepackage{hyperref}
\usepackage{booktabs}
\usepackage{xcolor}
\usepackage{enumitem}
\usepackage{cleveref}
\usepackage[backgroundcolor=blue!5, linecolor=blue!40, linewidth=1pt,
            innertopmargin=6pt, innerbottommargin=6pt,
            innerleftmargin=7pt, innerrightmargin=7pt]{mdframed}

\geometry{margin=2.2cm}

% ── Standard project macros ──────────────────────────────────────────────────
\newcommand{\bq}{\bm{q}}
\newcommand{\bp}{\bm{p}}
\newcommand{\bu}{\bm{u}}
\newcommand{\btau}{\bm{\tau}}
\newcommand{\bF}{\bm{F}}
\newcommand{\bJ}{\bm{J}}
\newcommand{\bM}{\bm{M}}
\newcommand{\bC}{\bm{C}}
\newcommand{\bG}{\bm{G}}

% ── Coloured assessment boxes ────────────────────────────────────────────────
\newenvironment{correct}{%
  \begin{mdframed}[backgroundcolor=green!8, linecolor=green!50!black]%
  \textbf{\color{green!50!black}\checkmark\; Correct.}\;}{\end{mdframed}}

\newenvironment{incorrect}{%
  \begin{mdframed}[backgroundcolor=red!8, linecolor=red!60!black]%
  \textbf{\color{red!60!black}$\boldsymbol{\times}$\; Needs correction.}\;}{\end{mdframed}}

\newenvironment{nuanced}{%
  \begin{mdframed}[backgroundcolor=yellow!10, linecolor=orange!60!black]%
  \textbf{\color{orange!60!black}$\sim$\; Partially correct / nuanced.}\;}{\end{mdframed}}

\title{\textbf{Exosuit Design Meeting Notes}\\
       \large Technical Analysis of Key Discussion Points\\[4pt]
       \normalsize \texttt{cable/cable\_with\_exo\_springs*.py}}
\date{April 2026}
\author{Isaac-sim robotics notes}

\begin{document}
\maketitle

% ═════════════════════════════════════════════════════════════════════════════
\section{Context}
% ═════════════════════════════════════════════════════════════════════════════

These notes document the design discussions around the exosuit
co-contraction mechanism for the cup manipulator.  Each point raised in the
meeting is stated, then assessed against the mathematical models derived in
the companion document
(\texttt{exosuit\_stiffness\_and\_co\_contraction.tex}).

\medskip\noindent
Two cable routing methods are under consideration:
\begin{itemize}[nosep]
  \item \textbf{Method~A} --- offset elbow pulleys
    ($r_{\text{elb}} = 32\;\text{mm}$, $d_{\text{off}} = 103\;\text{mm}$);
  \item \textbf{Method~B} --- centred elbow pulley
    ($r_{\text{cp}} = 30\;\text{mm}$ from mesh, on joint axis).
\end{itemize}


% ═════════════════════════════════════════════════════════════════════════════
\section{Point 1: Two Separate Pulleys vs.\ One Pulley with Two Grooves}
\label{sec:point1}
% ═════════════════════════════════════════════════════════════════════════════

\noindent
\textbf{Meeting claim:} ``Use two separate elbow pulleys, one for each exo
cable.''

\begin{nuanced}
This describes \textbf{Method~A}, which does use two separate elbow pulleys
(right and left, offset from the joint axis).  However, the current
Method~B implementation in
\texttt{cable\_with\_exo\_springs\_elbow\_follow.py} uses a
\textbf{single shared pulley} with \textbf{two grooves} separated by
$3\;\text{mm}$ (upper at $Z = 286.55\;\text{mm}$, lower at
$Z = 283.55\;\text{mm}$).

\medskip\noindent
Both approaches are valid but have different trade-offs:
\begin{itemize}[nosep]
  \item \textbf{Two separate pulleys} (Method~A): simpler tangent geometry
    (no groove-level Z tracking), but variable moment arm and
    $d_{\text{off}} \neq 0$.
  \item \textbf{Single pulley, two grooves} (Method~B): constant moment
    arm ($r_{\text{cp}}$), exact cable kinematics, but requires careful
    Z-plane management and encoder tracking.
\end{itemize}

\noindent
The code currently supports both.  The choice depends on whether the
prototype uses the centred design (Method~B) or offset mounting
(Method~A).
\end{nuanced}


% ═════════════════════════════════════════════════════════════════════════════
\section{Point 2: Stepper Motor Latency for Zero-Resistance Feel}
\label{sec:point2}
% ═════════════════════════════════════════════════════════════════════════════

\noindent
\textbf{Meeting claim:} ``A stepper motor can track the elbow encoder fast
enough to maintain zero cable tension during free rotation.''

\begin{nuanced}
The encoder-tracking law for Method~B requires the motor to feed cable at:
%
\begin{equation}
  \dot{l}_m = r_{\text{cp}} \cdot \dot{q}_2
\end{equation}
%
For a fast elbow motion $\dot{q}_2 = 5\;\text{rad/s}$ and
$r_{\text{cp}} = 47.75\;\text{mm}$:
%
\[
  \dot{l}_m = 0.04775 \times 5 = 0.239\;\text{m/s}
  = 239\;\text{mm/s}
\]

\noindent
\textbf{Latency analysis:}
A typical NEMA-17 stepper at~$3000\;\text{steps/s}$ with $1.8°$/step and
$20\;\text{mm}$ pulley gives
$v_{\max} \approx 300\;\text{mm/s}$ --- marginally sufficient.  However:

\begin{enumerate}[nosep]
  \item \textbf{Step-direction latency}: $\sim 1\;\text{ms}$ (acceptable).
  \item \textbf{Acceleration bandwidth}: Limited by rotor inertia and
    holding torque.  During sudden direction changes
    ($\ddot{q}_2 > 50\;\text{rad/s}^2$), the stepper may lose steps.
  \item \textbf{Microstepping resonance}: At mid-range speeds
    ($200$--$800\;\text{steps/s}$) steppers exhibit torque dips.
\end{enumerate}

\noindent
The spring extension due to tracking latency $\Delta t$:
\begin{equation}\label{eq:spring-lag}
  \delta_{\text{lag}} = r_{\text{cp}} \cdot \dot{q}_2 \cdot \Delta t
\end{equation}
For $\Delta t = 5\;\text{ms}$ and $\dot{q}_2 = 5\;\text{rad/s}$:
$\delta_{\text{lag}} = 0.04775 \times 5 \times 0.005 = 1.19\;\text{mm}$.
Parasitic force: $F = k_{\text{exo}} \times 0.00119 = 0.24\;\text{N}$
(with $k_{\text{exo}} = 200\;\text{N/m}$).

\medskip\noindent
\textbf{Verdict:} At slow motions ($\dot{q}_2 < 2\;\text{rad/s}$), stepper
tracking is adequate.  At fast motions with sudden reversals, the parasitic
force may be perceptible.  A \textbf{servo motor with encoder feedback}
(e.g., closed-loop stepper or small BLDC) would provide tighter tracking.
\end{nuanced}


% ═════════════════════════════════════════════════════════════════════════════
\section{Point 3: CubeMars Impedance Control Eliminating Compliance}
\label{sec:point3}
% ═════════════════════════════════════════════════════════════════════════════

\noindent
\textbf{Meeting claim:} ``Using a CubeMars motor in impedance/torque mode
could eliminate the need for physical springs entirely --- the motor itself
provides the compliance.''

\begin{nuanced}
This is conceptually attractive but has important limitations:

\medskip\noindent
\textbf{What CubeMars MIT torque mode provides:}
\begin{itemize}[nosep]
  \item Direct torque control at $\sim 1\;\text{kHz}$ loop rate
  \item Programmable virtual stiffness: $\tau = K_p(\theta_{\text{des}} - \theta) + K_d(\dot\theta_{\text{des}} - \dot\theta)$
  \item AK60-6: $\tau_{\text{peak}} = 9\;\text{N\,m}$,
    $J_m = 3.32 \times 10^{-5}\;\text{kg\,m}^2$, $N = 6$
\end{itemize}

\medskip\noindent
\textbf{Why physical springs are still valuable:}
\begin{enumerate}[nosep]
  \item \textbf{Impact bandwidth:} An impact at the elbow has a
    characteristic time $\sim 1$--$5\;\text{ms}$.  The CubeMars control
    loop runs at $\sim 1\;\text{ms}$.  A physical spring responds at the
    \emph{speed of sound in the material} ($\ll 0.1\;\text{ms}$) ---
    orders of magnitude faster.
  \item \textbf{Energy storage:} Springs absorb impact energy passively.
    A motor must dissipate energy electrically (back-EMF, resistive losses).
  \item \textbf{Fail-safe:} If motor power is lost, a pre-tensioned spring
    still provides stiffness.  A virtual spring provides zero stiffness
    when powered off.
  \item \textbf{Motor resonance:} The reflected rotor inertia
    $J_{\text{eff}} = N^2 \cdot J_m = 36 \times 3.32 \times 10^{-5}
    = 1.2 \times 10^{-3}\;\text{kg\,m}^2$ creates a resonance with the
    virtual spring:
    \begin{equation}
      \omega_n = \sqrt{\frac{K_p}{J_{\text{eff}}}}
    \end{equation}
    For $K_p = 10\;\text{N\,m/rad}$:
    $\omega_n = \sqrt{10 / 0.0012} \approx 91\;\text{rad/s}$
    ($14.5\;\text{Hz}$).  Above this frequency, the motor cannot track.
\end{enumerate}

\medskip\noindent
\textbf{Verdict:} CubeMars impedance mode is excellent for \emph{low-frequency}
stiffness modulation (trajectory following, gravity compensation) but
\textbf{cannot replace physical springs for impact protection}.  The
recommended architecture is \textbf{hybrid}: physical springs for
high-bandwidth passive compliance + CubeMars for active stiffness tuning.
This is exactly the SEA (Series Elastic Actuator) design already implemented
in \texttt{actuators/sea.py}.
\end{nuanced}


% ═════════════════════════════════════════════════════════════════════════════
\section{Point 4: The ``$2k$'' Co-Contraction Stiffness}
\label{sec:point4}
% ═════════════════════════════════════════════════════════════════════════════

\noindent
\textbf{Meeting claim:} ``When both exo cables co-contract, the effective
joint stiffness is $2k$ times some constant.''

\begin{correct}
This is exactly the co-contraction result derived in the companion document.
For symmetric pretension on antagonistic cables, the effective stiffness at
the joint is:
%
\begin{equation}\label{eq:keff-general}
  \boxed{k_{\text{eff}} = 2 \, k_{\text{exo}} \, r_{\text{eff}}^2}
\end{equation}
%
where the factor of~$2$ arises because \emph{both} cables contribute to
restoring torque (one cable gets longer, the other shorter, and both resist
the deflection).  The effective moment arm $r_{\text{eff}}$ depends on the
method:

\begin{itemize}[nosep]
  \item \textbf{Method~A:}
    $r_{\text{eff}} = l_c'(q_2)
    \approx r_{\text{elb}} = 32\;\text{mm}$
    $\;\Rightarrow\;
    k_{\text{eff}}^A \approx 2 \times 200 \times 0.032^2
    = 0.41\;\text{N\,m/rad}$

  \item \textbf{Method~B:}
    $r_{\text{eff}} = r_{\text{cp}} = 47.75\;\text{mm}$
    $\;\Rightarrow\;
    k_{\text{eff}}^B = 2 \times 200 \times 0.04775^2
    = 0.91\;\text{N\,m/rad}$
\end{itemize}

\noindent
\textbf{Derivation sketch:}
A deflection $\Delta q$ from equilibrium changes each cable's spring
extension by $\pm r_{\text{eff}} \cdot \Delta q$.  Each spring contributes
force $k_{\text{exo}} \cdot r_{\text{eff}} \cdot \Delta q$, and each force
acts through moment arm $r_{\text{eff}}$, giving torque
$k_{\text{exo}} \cdot r_{\text{eff}}^2 \cdot \Delta q$ per cable.  Two
cables $\Rightarrow$ factor of~2.
\end{correct}


% ═════════════════════════════════════════════════════════════════════════════
\section{Point 5: Springs on the Upper Arm Side}
\label{sec:point5}
% ═════════════════════════════════════════════════════════════════════════════

\noindent
\textbf{Meeting claim:} ``Place the springs on the upper arm (link~1) side
of the elbow pulley.''

\begin{incorrect}
In the current architecture, the springs are placed \textbf{between the
elbow pulley and the link-2 anchor} --- i.e., on the \textbf{forearm}
(link~2) side, not the upper arm (link~1) side.

\medskip\noindent
\textbf{Why the forearm side is correct:}
\begin{enumerate}[nosep]
  \item The cable segment from motor $\to$ L1 pulley $\to$ elbow pulley must
    be \emph{inextensible} --- the motor must control cable displacement
    precisely.  A spring on this side would absorb motor motion and destroy
    position authority.
  \item The spring on the link-2 side acts as a \emph{force sensor}: spring
    extension $\delta$ directly indicates cable tension
    $F = k_{\text{exo}} \cdot \delta$.
  \item From the TikZ architecture diagram
    (\texttt{exosuit\_stiffness\_and\_co\_contraction.tex},
    \cref{fig:architecture}): the cable path is
    Motor $\to$ cable $\to$ L1 pulley $\to$ cable $\to$ Elbow pulley
    $\to$ \textbf{spring} $\to$ cable $\to$ End anchor (link~2).
  \item The spring provides compliance at the \emph{impact site} (link~2
    interacts with objects), which is the correct place for impact
    absorption.  A spring on the motor side would delay force transmission
    rather than absorb impacts locally.
\end{enumerate}

\medskip\noindent
In the code (\texttt{cable\_with\_exo\_springs\_elbow\_follow.py}), the
spring visual is drawn as a fraction (\texttt{spring\_fraction} $= 0.12$)
of the cable segment between the elbow pulley and the end anchor --- i.e.,
on the link-2 side.
\end{incorrect}


% ═════════════════════════════════════════════════════════════════════════════
\section{Point 6: Transmission Ratio via Pulley Diameters}
\label{sec:point6}
% ═════════════════════════════════════════════════════════════════════════════

\noindent
\textbf{Meeting claim:} ``Changing the motor pulley diameter relative to the
elbow pulley changes the transmission ratio and thus the stiffness.''

\begin{correct}
The transmission ratio between motor and joint is determined by the pulley
radii.  Define:
%
\begin{align}
  r_m &= \text{motor pulley radius}, \\
  r_{\text{elb}} &= \text{elbow pulley radius (or } r_{\text{cp}} \text{ for Method~B)}.
\end{align}
%
The kinematic relation is:
\begin{equation}
  l_m = r_m \cdot \theta_m, \qquad
  l_c = r_{\text{elb}} \cdot q_2
  \qquad\Rightarrow\qquad
  \frac{\theta_m}{q_2} = \frac{r_{\text{elb}}}{r_m}
\end{equation}
%
If the motor and elbow pulleys are the same radius
($r_m = r_{\text{elb}} = r_{\text{cp}}$), the ratio is 1:1.

\medskip\noindent
\textbf{Effect on stiffness:}
The spring extension becomes:
\begin{equation}
  \delta = r_m \cdot \theta_m - r_{\text{elb}} \cdot q_2
\end{equation}
The effective joint stiffness with co-contraction:
\begin{equation}
  k_{\text{eff}} = 2 \, k_{\text{exo}} \, r_{\text{elb}}^2
\end{equation}
Note: the elbow-side radius determines joint stiffness (it is the output
lever arm).  Changing $r_m$ affects the \emph{motor-side resolution and
torque requirement}:
\begin{itemize}[nosep]
  \item \textbf{Smaller $r_m$}: motor sees higher mechanical advantage,
    needs less torque but more rotation per unit cable displacement.
  \item \textbf{Larger $r_m$}: motor sees lower mechanical advantage,
    needs more torque but less rotation.
\end{itemize}

\noindent
Currently, Method~B uses $r_m = r_{\text{cp}} = 47.75\;\text{mm}$ (1:1
ratio), which simplifies the encoder tracking law to
$\theta_m = q_2 + \Delta\theta$.
\end{correct}


% ═════════════════════════════════════════════════════════════════════════════
\section{Point 7: Encoder Tracking Necessity for Method~B}
\label{sec:point7}
% ═════════════════════════════════════════════════════════════════════════════

\noindent
\textbf{Meeting claim:} ``Method~B (centred pulley) requires continuous
encoder tracking to prevent parasitic forces during free rotation.''

\begin{correct}
This is the fundamental requirement of Method~B.  Because the elbow pulley
is \emph{on} the joint axis, any rotation $\Delta q_2$ wraps/unwraps cable
by exactly $r_{\text{cp}} \cdot \Delta q_2$.  Without motor compensation,
the spring extends:
%
\begin{equation}
  \delta = r_{\text{cp}} \cdot \Delta q_2
\end{equation}
%
producing a parasitic restoring torque:
\begin{equation}
  \tau_{\text{parasitic}}
  = k_{\text{exo}} \cdot r_{\text{cp}}^2 \cdot \Delta q_2
  = k_{\text{eff}}^B \cdot \Delta q_2
\end{equation}

For $\Delta q_2 = 10°\;(0.175\;\text{rad})$:
$\tau_{\text{parasitic}} = 0.91 \times 0.175 = 0.16\;\text{N\,m}$
--- clearly perceptible and unacceptable for ``transparent'' free rotation.

\medskip\noindent
The tracking law is:
\begin{equation}
  \theta_{m,R} = q_2 + \Delta\theta, \qquad
  \theta_{m,L} = -q_2 + \Delta\theta
\end{equation}
where $\Delta\theta = 0$ for exo~OFF (zero force) and
$\Delta\theta > 0$ for co-contraction (symmetric pretension).

\medskip\noindent
\textbf{Contrast with Method~A:}
In Method~A, the springs sit at rest \emph{passively} because the offset
geometry does not couple $q_2$ rotation to spring extension (only the motor
side of the cable connects to the offset pulley, and the spring is only
loaded when the motor actively pulls).  The encoder is an \emph{upgrade},
not a requirement.
\end{correct}


% ═════════════════════════════════════════════════════════════════════════════
\section{Summary of Assessments}
\label{sec:summary}
% ═════════════════════════════════════════════════════════════════════════════

\begin{table}[h]
\centering
\small
\begin{tabular}{@{}clc@{}}
\toprule
\textbf{\#} & \textbf{Discussion Point} & \textbf{Assessment} \\
\midrule
1 & Two separate pulleys vs.\ one with two grooves
  & \textcolor{orange!60!black}{Nuanced} \\
2 & Stepper latency for zero-resistance feel
  & \textcolor{orange!60!black}{Nuanced} \\
3 & CubeMars impedance eliminating springs
  & \textcolor{orange!60!black}{Nuanced} \\
4 & ``$2k$'' co-contraction stiffness
  & \textcolor{green!50!black}{Correct} \\
5 & Springs on the upper arm side
  & \textcolor{red!60!black}{Incorrect} \\
6 & Transmission ratio via pulley diameters
  & \textcolor{green!50!black}{Correct} \\
7 & Encoder tracking required for Method~B
  & \textcolor{green!50!black}{Correct} \\
\bottomrule
\end{tabular}
\caption{Summary of meeting point assessments.}
\label{tab:summary}
\end{table}


% ═════════════════════════════════════════════════════════════════════════════
\section{Action Items}
\label{sec:actions}
% ═════════════════════════════════════════════════════════════════════════════

\begin{enumerate}
  \item \textbf{Prototype decision:} Choose Method~A or~B for the first
    physical build.  Method~B is mechanically simpler (one pulley, exact
    kinematics) but requires reliable encoder tracking.

  \item \textbf{Stepper vs.\ servo motor test:} Build a bench test with a
    stepper (or closed-loop stepper) tracking an encoder at
    $\dot{q}_2 \leq 5\;\text{rad/s}$.  Measure tracking latency and
    parasitic spring force.

  \item \textbf{Spring placement:} Confirm springs are between elbow pulley
    exit and link-2 anchor (forearm side), not on the motor/upper-arm side.

  \item \textbf{CubeMars integration path:} If CubeMars motors are used for
    the exo cables, implement hybrid control: physical springs for
    high-bandwidth impact response + motor torque mode for active stiffness
    modulation.

  \item \textbf{Pulley diameter trade study:} Evaluate whether a non-unity
    transmission ratio ($r_m \neq r_{\text{cp}}$) improves motor torque
    headroom or resolution.

  \item \textbf{Simulation validation:} Add exo cable dynamics to the
    Drake or Isaac~Sim pipeline.  Use the SEA infrastructure
    (\texttt{actuators/sea.py}) as a template for the exo cable
    spring-damper model.
\end{enumerate}


% ═════════════════════════════════════════════════════════════════════════════
\section{Progress Update --- 21 April 2026}
\label{sec:progress-apr21}
% ═════════════════════════════════════════════════════════════════════════════

\noindent
\textit{Recorded meeting: \texttt{Audio-Weekly progress-20260421\_115439-Meeting Recording.mp4}.
Attendees: Debarup (in-person), Chris (remote, in lab), Ben (joined mid-meeting).
Tube/rail strike restricted in-person attendance.}

% ─────────────────────────────────────────────────────────────────────────────
\subsection{Arduino Motor Driver}
% ─────────────────────────────────────────────────────────────────────────────

\begin{itemize}[nosep]
  \item The wrong driver issue has been \textbf{resolved}.
  \item Two candidate Arduino libraries found for communicating with the
    DeBio motor controller via Arduino Uno:
    \begin{enumerate}[nosep]
      \item \textbf{Library~1} (PlatformIO-compatible, no English docs) ---
        easy to install via PlatformIO but documentation is in Chinese only.
      \item \textbf{Library~2} (English docs, not on PlatformIO) ---
        documentation available but requires manual installation.
    \end{enumerate}
  \item Most library function names are in English; inferred usage despite
    Chinese documentation.  Example mistranslation: ``point'' rendered as
    ``indicator''.
  \item \textbf{Recommended translation tool:} \textbf{DeepL}
    (\url{https://www.deepl.com}) for technical documentation.
  \item \textbf{Recommended workflow:} Clone the library's GitHub repo,
    open in VS~Code, and use GitHub~Copilot to auto-generate English
    documentation from the English source code.
\end{itemize}

% ─────────────────────────────────────────────────────────────────────────────
\subsection{Physical Hardware --- Cable Stoppers and Inserts}
% ─────────────────────────────────────────────────────────────────────────────

\begin{itemize}[nosep]
  \item Currently \textbf{3D-printing replacement parts} for the cable
    infrastructure (inserts and loop stoppers).
  \item Need a \textbf{cable-loop stopper} (white stopper previously used for
    the wire loop).  Attempting to replicate with 3D-printed plastic;
    if insufficient, will order metal stoppers.
  \item Additional bolts and inserts may also need ordering; list to be sent
    by EOD 21~April.
  \item \textbf{Power supply issue:} Desktop power supply has non-standard
    plugs; may swap for a screw-terminal unit to connect motor and test rig.
\end{itemize}

% ─────────────────────────────────────────────────────────────────────────────
\subsection{Stepper vs.\ Servo Motor Decision}
\label{sec:motor-choice}
% ─────────────────────────────────────────────────────────────────────────────

\begin{itemize}[nosep]
  \item Original plan was to use a \textbf{stepper motor} to track the joint
    encoder and maintain near-zero cable tension.
  \item Key concern: steppers are \textbf{not back-drivable} and can
    \textbf{miss steps} at high speed or under load---this adds perceptible
    resistance when moving the exo freely.
  \item \textbf{Servo motor alternative:} Peiyu (lab colleague) provided a
    servo motor (brought from China for personal use).  This will be the
    first hardware under test.
  \item If the servo hypothesis validates, cheap but decent servo motors will
    be ordered from Robot Shop (fast delivery).  The CubeMars motors are
    reserved for the manipulator, not the exo cables.
  \item The stepper latency test plan remains valid but the servo is now
    the \textbf{primary candidate}.
\end{itemize}

% ─────────────────────────────────────────────────────────────────────────────
\subsection{CubeMars Motor Procurement}
\label{sec:cubemarks-procurement}
% ─────────────────────────────────────────────────────────────────────────────

\begin{itemize}[nosep]
  \item \textbf{Compliance certificate approved.}
  \item Two CubeMars motors ordered for the \textbf{manipulator} (not exo).
  \item Procurement path:
    \begin{enumerate}[nosep]
      \item 3DXR registered as a \textbf{university approved supplier} (ICT
        request submitted 21~April).
      \item HALA asked to create a \textbf{VCC (Virtual Credit Card)} for
        the order.  VCC approval typically takes 4--5 days.
      \item Once VCC is issued, order placed; stock is available and
        delivery is $\approx$1--2 days.
    \end{enumerate}
  \item Motors feature a \textbf{planetary gearbox} (ratio $\approx 1:6$,
    back-drivable).  These are the same class of motor used in
    humanoid/compliant robots.
  \item Exo cable motors remain a separate problem requiring the
    servo/stepper latency test.
\end{itemize}

% ─────────────────────────────────────────────────────────────────────────────
\subsection{GitHub Copilot Student Access}
% ─────────────────────────────────────────────────────────────────────────────

\begin{itemize}[nosep]
  \item Free access via the \textbf{GitHub Student Developer Pack}.
    Chris already has this.  Strongly recommended for translating Chinese
    library docs and generating docstrings.
\end{itemize}

% ─────────────────────────────────────────────────────────────────────────────
\subsection{Lab Access --- Robotics Area}
% ─────────────────────────────────────────────────────────────────────────────

\begin{itemize}[nosep]
  \item Ben (Chris?) submitted a request for \textbf{general robotics lab
    access} on 21~April.  Parker (lab manager) to register them.
  \item General induction already completed.  Access list update pending
    Becky's confirmation.
  \item Currently using the Ideas Lab upstairs mainly for 3D printer access;
    robotics area useful for hardware testing.
\end{itemize}

% ─────────────────────────────────────────────────────────────────────────────
\subsection{Exosuit Architecture --- Key Design Discussion (with Ben)}
\label{sec:exo-architecture}
% ─────────────────────────────────────────────────────────────────────────────

Ben joined mid-meeting and described the wearable exosuit hardware.  A major
architectural discussion followed on how the exosuit, manipulator, and
contraction-theory controller should be integrated.

\subsubsection*{Ben's Wearable Exosuit Hardware}

\begin{itemize}[nosep]
  \item \textbf{Current wearable prototype:} single fixed elbow pulley +
    two separate braided cables routed to small motors.  The motors can
    independently tension the exo springs.
  \item The two motors (``purple'' cables in the diagram) \emph{do not drive
    the joint} --- they only add stiffness.  The joint moves freely; the
    cables create a symmetric restoring force (co-contraction stiffness
    $k_{\text{eff}} = 2 k_{\text{exo}} r_{\text{eff}}^2$).
  \item The adapter pulley is \textbf{mounted on the joint shaft}, so it
    rotates with the joint.  Without motor compensation, this wraps/unwraps
    cable and stretches the springs (parasitic stiffness, see
    \cref{sec:point7}).  \textbf{Active encoder tracking nullifies this.}
\end{itemize}

\subsubsection*{Exosuit as a Signal Source (Not a Force Actuator)}

The central design insight agreed in this meeting:

\begin{mdframed}[backgroundcolor=blue!5, linecolor=blue!40]
\textbf{Key principle:} The exosuit is \emph{not} meant to generate the
required stiffness directly.  It acts as a \textbf{stiffness signal source}
that tells the manipulator/robot how much stiffness to generate.  The
manipulator's own back-drivable motors implement that stiffness via
co-contraction.
\end{mdframed}

\noindent
Consequences:
\begin{itemize}[nosep]
  \item The exosuit motors need only be small and lightweight --- they give
    a $\delta\theta$ command (tug), not a full torque command.
  \item When the manipulator already has the correct stiffness, the exosuit
    commands $\delta\theta = 0$ (stays passive).
  \item When a disturbance arrives and stiffness is insufficient, the exosuit
    activates: applies $q_2 \pm \delta\theta$ to stiffen the spring in the
    opposing direction.  This is exactly the $\Delta\theta$ co-contraction
    command already implemented in simulation.
  \item Analogy used in meeting: ``like a horse rider --- you keep giving
    signals until the horse picks up speed, then you stop.''
\end{itemize}

\subsubsection*{Three-Layer Control Architecture}

The full system consists of three independent but cooperating layers:

\begin{enumerate}
  \item \textbf{Motion control} (manipulator joint trajectories) ---
    computed-torque / LQR controller.
  \item \textbf{Stiffness control} (co-contraction via manipulator exo
    springs + back-drivable motors) --- driven by the contraction-theory
    algorithm (\texttt{contraction-theory/}).  This layer decides the
    correct $k_{\text{eff}}$ given the disturbance.
  \item \textbf{Exosuit signal layer} (wearable, lightweight) --- a third
    add-on that monitors stiffness error and sends $\delta\theta$ signals to
    Layer~2 to ``teach'' it the correct stiffness level.
\end{enumerate}

\noindent
\textbf{Impairment analogy:} Treat the manipulator's stiffness controller
as an ``impaired brain'' that does not know the optimal co-contraction level.
The exosuit (Layer~3) is the rehabilitation tool: it provides an external
stiffness signal until Layer~2 learns the correct response.  Once Layer~2
converges, Layer~3 becomes passive.

\subsubsection*{Are Manipulator Springs Still Needed?}

\begin{itemize}[nosep]
  \item If the manipulator uses \textbf{back-drivable motors} (e.g.,
    CubeMars), physical springs in the cable path may be redundant ---
    the motor itself provides programmable stiffness.
  \item Springs are still needed for the \textbf{exo cable motors} if those
    use steppers (not back-drivable).
  \item Conclusion: \textbf{springs may be removed from the manipulator
    cable path} once CubeMars motors are installed, but retained on the
    exo cable motors until a back-drivable alternative is confirmed.
\end{itemize}

\subsubsection*{Stiffness Modulation Ideas Discussed}

Two lightweight alternatives to continuous motor pretension were discussed,
motivated by the desire to make the wearable exosuit as small and
power-efficient as possible.

% ── PWM ─────────────────────────────────────────────────────────────────────
\paragraph{Option A: PWM stiffness modulation.}

The motor drives the cable at full amplitude but is switched on/off at a
high duty cycle $D \in [0,1]$.  When ON the spring is engaged at its full
stiffness $k_{\text{exo}}$; when OFF the motor back-drives freely and the
spring goes slack.  Averaged over one PWM period, the effective joint
stiffness is approximately:
%
\begin{equation}\label{eq:pwm-keff}
  k_{\text{eff,PWM}} \approx D \cdot 2\,k_{\text{exo}}\,r_{\text{cp}}^2
\end{equation}
%
so varying $D$ from 0 to 1 continuously sweeps $k_{\text{eff}}$ from zero
to its maximum value.

\medskip\noindent
\textbf{Condition for validity:}  The PWM frequency $f_{\text{PWM}}$ must be
much higher than the mechanical bandwidth of the joint:
%
\begin{equation}
  f_{\text{PWM}} \gg \frac{1}{2\pi}\sqrt{\frac{k_{\text{eff,max}}}{J_{\text{link}}}}
  \approx \frac{1}{2\pi}\sqrt{\frac{0.91}{0.05}} \approx 21\;\text{Hz}
\end{equation}
%
so $f_{\text{PWM}} > 300\;\text{Hz}$ (well within electrical limits).
However, the \emph{mechanical} system (motor inertia + cable compliance)
may not follow at that frequency, causing the joint to feel discrete
on/off pulses as jitter rather than smooth intermediate stiffness.
\textbf{Needs experimental validation.}

% ── Clutch ───────────────────────────────────────────────────────────────────
\paragraph{Option B: Clutch / brake-pad (binary stiffness switch).}

Replace the motor+spring with a \textbf{lockable cable mechanism} --- a
spring pre-tensioned at assembly, held free by a clutch.  On command
(solenoid, shape-memory wire, or micro-servo), the clutch engages and locks
the cable, instantly loading the spring.

\medskip\noindent
The torque applied to the joint becomes:
%
\begin{equation}\label{eq:clutch-tau}
  \tau_{\text{exo}} =
  \begin{cases}
    k_{\text{exo}}\,r_{\text{cp}}\,\bigl(r_{\text{cp}}\,\delta q - \delta_0\bigr)
      & \text{clutch ON (spring engaged)} \\
    0 & \text{clutch OFF (cable free)}
  \end{cases}
\end{equation}
%
where $\delta_0$ is the spring pre-compression at rest and $\delta q$ is
the joint deflection from nominal.

\medskip\noindent
\textbf{Effective behaviour:}
\begin{itemize}[nosep]
  \item Engagement is \emph{mechanical} speed --- not limited by a
    control loop.  Response is sub-millisecond for solenoid actuation.
  \item No continuous power is required to hold stiffness (spring is
    passive once locked).
  \item The energy absorbed when the clutch engages mid-motion goes into
    the spring (elastic storage), not heat --- safe for wearable use.
  \item The on/off nature provides \emph{bang-bang stiffness}: either full
    $k_{\text{eff}}$ or zero.  For the signal-source use case this is
    sufficient --- the exosuit only needs to indicate ``stiffen now'',
    not modulate a precise level.
\end{itemize}

\medskip\noindent
\textbf{Comparison:}

\begin{table}[ht]
\centering\small
\begin{tabular}{@{}lccccc@{}}
\toprule
\textbf{Mechanism} & \textbf{Stiffness} & \textbf{Complexity}
  & \textbf{Weight} & \textbf{Power (hold)} & \textbf{Bandwidth} \\
\midrule
Motor + spring (current) & Continuous & High & Medium & High & Servo BW \\
PWM motor               & Pseudo-variable & Medium & Medium & Medium & Needs test \\
Clutch / brake-pad      & Binary ($0$ or $k_s$) & Low & Low & None & Mechanical \\
\bottomrule
\end{tabular}
\caption{Comparison of exosuit stiffness modulation options.}
\label{tab:stiffness-options}
\end{table}

\noindent
For a \textbf{lightweight wearable signal source}, the clutch option is the
most attractive: a small solenoid (mass $<30\;\text{g}$) or shape-memory
actuator can lock/unlock the cable, requiring only a brief current pulse
(not sustained power) to change state.

% ─────────────────────────────────────────────────────────────────────────────
\subsection{Agreed Next Steps}
% ─────────────────────────────────────────────────────────────────────────────

\begin{enumerate}
  \item \textbf{[Chris, \(\leq\)28~Apr]} Complete stepper/servo latency
    test: measure tracking lag between joint encoder and exo motor at
    $\dot{q}_2 \leq 2\;\text{rad/s}$.  Use Peiyu's servo motor as first
    candidate.  If lag is acceptable, order cheap servo motors from Robot
    Shop.
  \item \textbf{[Chris, EOD 21~Apr]} Confirm 3D-printed cable-loop stopper
    feasibility; send purchase list for stoppers and bolts if 3D printing
    inadequate.
  \item \textbf{[Chris]} Resolve Arduino library choice; use DeepL /
    Copilot for Chinese documentation.
  \item \textbf{[Chris]} Register for GitHub Student Developer Pack.
  \item \textbf{[Debarup]} Finalise VCC with HALA; order two CubeMars
    motors from 3DXR once ICT approves supplier.
  \item \textbf{[Debarup/Ben]} Implement contraction-theory stiffness
    controller on the manipulator \emph{first} (without exosuit), validate
    that it generates correct $k_{\text{eff}}$ under disturbances.
  \item \textbf{[All]} Add the wearable exosuit signal layer as a third
    add-on once Layers~1 and~2 are stable.
  \item \textbf{[Ben]} Review CubeMars motor spec (planetary gear, 1:6
    ratio, back-drivable) against exo motor requirements via links already
    shared.
\end{enumerate}

\end{document}
